% DECLARATION OF INDIVIDUAL CONTRIBUTIONS
% ======================================================================
\newpage
\begin{center}
{\fontsize{16}{19}\selectfont\textbf{Declaration of Individual Contributions}}
\end{center}

\vspace{3mm}

In accordance with the University of Hong Kong's guidelines for group project reports, this section details the specific contributions of each team member.

\vspace{4mm}

\noindent\textbf{Lin Chuanbin:}

\vspace{2mm}

In the Agent Platform project, I was responsible for developing three core modules from scratch: the Tool Agent microservice, the User Service microservice, and the full-stack implementation of the initial front-end version, delivering a total of 30+ core files. Below are my specific responsibilities:

\textbf{1 Tool Agent Module}

The Tool Agent is the core intelligent agent of the system, built on Spring Boot, integrating the DeepSeek large language model to achieve natural language understanding and intent recognition. It covers five subsystems: AI intent routing, meeting room management, schedule conflict detection, intelligent route planning, and WebSocket real-time communication.

1.1 System Architecture and AI Intent Recognition

I designed a layered architecture for the Tool Agent, clearly defining the Controller layer, Service layer, external service layer, and data access layer. The Controller layer provides a unified RESTful API entry point; the Service layer is responsible for core business functions such as intent routing, meeting room management, and schedule conflict detection; the external service layer encapsulates DeepSeek AI calls and the Gaode Map API; and the data layer implements database access through MyBatis-Plus.

The AI intent recognition module is the intelligent brain of the Tool Agent. I encapsulated the communication logic with the DeepSeek API, using a RESTful approach and supporting a dual-role dialogue mode of system prompts and user content. The system defines three intent types (route planning, meeting room search, and schedule conflict detection). After DeepSeek analyzes the intent, it automatically routes the request to the corresponding business processing method.

Considering the potential unavailability of AI services in a production environment, I implemented a robust fallback mechanism: when a DeepSeek call fails, the system returns preset demo data based on the tool type, ensuring the front-end always displays correctly and the user experience is not affected by AI service failures.

1.2 Tool Agent Module

The meeting room module contains two core data models: meeting room entities and schedule booking entities. I implemented time period overlap detection logic, returning complete data including availability status and a list of devices. The booking function first checks for conflicts; if the time period is already booked, it returns a failure.

For AI intelligent matching, the system calls DeepSeek to analyze the user's meeting room needs, extracting information such as date, number of people, and equipment. The returned results include an AI matching score and recommendation reasons, helping users quickly find the most suitable meeting room.

1.3 Schedule Conflict Detection (High-Performance Redis Implementation)

Schedule data is stored in a Redis ordered set, using user ID as the key and timestamp as the score, achieving range queries with O(log N) time complexity. I designed three core interfaces: range query conflict detection, schedule data storage, and record deletion by event ID.

Timestamp calculation uses the East 8 time zone to ensure consistency with domestic business time.

1.4 Intelligent Route Planning (Gaode API Three-Level Degradation)

I encapsulated the calls to the Gaode Map REST API, mainly using three interfaces: geocoding, POI search, and driving route planning. To address the address resolution success rate issue, I designed a three-level degradation strategy: first, attempt direct geocoding; if it fails and the address does not contain a city prefix, automatically append the "Beijing" prefix and retry (this was later modified); finally, use POI text search as a fallback, significantly improving the address resolution success rate.

The system formats the raw data returned by Gaode Maps in a user-friendly way: distance is automatically converted to "km/m" units, duration to "hours/minutes" units, and traffic conditions are automatically classified into three levels: "smooth/slow/congested" based on average speed.

1.5 WebSocket Real-time Task Progress Push

The system provides a real-time task progress push service via WebSocket, supporting the establishment of independent connections by task ID. A concurrent and secure session management mechanism is used to manage all active connections, simulating a complete processing flow (parsing natural language → querying the database → AI model inference → organizing results), with progress updates pushed at fixed intervals for each stage. A thread pool is used to manage task execution, supporting a maximum of 10 concurrent tasks.

\textbf{2 User Service Module}

The User Service is the system's user authentication and permission management microservice, built on Spring Boot + Spring Security, using JWT to implement stateless identity authentication, supporting encrypted password storage and interface-level permission control. Yao Jiacheng and Xia Xin subsequently made some improvements and developments on this basis.

2.1 System Architecture Design

The User Service adopts a layered architecture, including a Controller layer, a Service layer, a data access layer, and an entity/DTO layer. Global security control is implemented through Spring Security configuration, including CSRF disabling, stateless session management, allowing login interfaces, and authentication protection for other interfaces.

2.2 JWT Authentication and Security Module

I encapsulated complete lifecycle management of JWTs, including token generation, parsing, and verification. The HMAC-SHA256 signature algorithm is used, with the key and expiration time injected from the configuration file for flexible configuration management. The generation method encodes the username and role information into the token; the parsing method verifies the token signature; the verification method catches all exceptions and returns a boolean value to ensure the security of the caller.

The security policy is configured with complete Spring Security protection: disabling CSRF (using JWT instead of Session in a front-end/back-end separation architecture), setting up stateless session management, allowing login interfaces, and requiring authentication for all other requests. Passwords are stored using BCrypt with one-way hash encryption.

2.3 User Login and Permission Management

The login process implements a complete four-step security verification: querying users by username, checking user status, BCrypt password matching verification, generating a token, and encapsulating the response for return.

I designed standardized request and response DTOs, using annotations to implement parameter validation to ensure that usernames and passwords are not empty. I also defined a unified generic response encapsulation class containing three fields: status code, message, and data, providing static factory methods for success and failure to ensure consistent return formats across all microservices.

\textbf{3 Front-end Development (Initial Version)}

I was responsible for building the initial version of the front-end project and implementing the core pages, providing the basic framework for subsequent iterative development by team members (later, Yao Jiacheng and Xia Xin proposed many excellent new ideas and made many significant improvements). My front-end uses the Vue 3 + TypeScript + Vite technology stack, combined with the Element Plus UI component library and Pinia state management.

3.1 Front-end Technology Architecture and Engineering

The project uses Vue 3's composable API to organize code logic, TypeScript for type safety, and Vite as the build tool for rapid development and hot updates. I configured path aliases to significantly improve code maintainability. I also configured API proxies and WebSocket proxies to solve cross-domain issues in the development environment. Element Plus uses on-demand plugin loading to import components as needed, reducing the package size.

3.2 Core Page Implementation (Initial Version)

Login Page: Uses a centered card layout with a gradient background for a modern feel. A form component is used for username and password input, with validation rules ensuring input integrity. The login logic calls a state management method; upon successful login, the token is persisted and the user is redirected to the homepage.

Main Layout Page: Uses a classic sidebar + top navigation + main content area layout. The sidebar implements a multi-level menu, and the top navigation displays current user information and logout functionality. The main content area uses transition animations to switch pages.

The tool call page adopts a left-right split layout. The left side is the function operation area (three tabs: meeting room search, schedule conflict detection, and route planning), and the right side is the AI assistant panel (natural language input, WebSocket real-time progress, and AI analysis result display). The map component encapsulates the initialization, route drawing, origin and destination marking, and view adaptation functions of Gaode Map, and uses a dynamic script loading method to import the map API.

Note: The above front-end components and pages are the initial version. Subsequent versions were refactored and expanded by Yao Jiacheng and Xia Xin, including a multi-tab system, responsive mobile adaptation, a global Copilot component, and a unified UI design system—all production-grade improvements.

3.3 State Management and HTTP Encapsulation

I designed a user state management module, defined using Pinia's Setup Store pattern, including responsive states such as tokens and user information, as well as Action methods such as login, logout, and retrieving user information. Upon successful login, the token is persisted to local storage; upon logout, the state is cleared and the user is redirected back to the login page. A task state management module was also designed, supporting adding tasks and updating task states.

The HTTP request encapsulation uses Axios to implement a unified request instance, configured with a relatively long timeout (to adapt to AI inference time), a basic path, and JSON content type. The request interceptor automatically reads the token and adds it to the request header. The response interceptor handles backend responses uniformly: returning data on success and displaying an error message on failure; unauthorized status triggers login expiration confirmation, and the user is automatically logged out and redirected to the login page after confirmation.

3.4 WebSocket Client and Task Orchestration (Initial Version)

I encapsulated a WebSocket client class, providing core methods such as connect, send, close, and event listeners. It supports an automatic reconnection mechanism (up to 3 times, 2-second interval), and achieves loosely coupled message processing through an event listener pattern. A singleton pattern is used to ensure global sharing of the same connection.

Simultaneously, I wrote the initial gateway center orchestration logic, implementing static Agent routing and WebSocket progress updates, providing the infrastructure for the subsequent implementation of asynchronous execution pipelines, failover, and retry mechanisms.

\vspace{4mm}

\noindent\textbf{Yao Jiacheng:}

\vspace{2mm}

I served as the backend architecture lead and, over the course of the project, took on full-stack development, system integration, security, and DevOps. Much of my work involved extending the initial implementations contributed by other members into production-grade features.

\textbf{1. Backend}

1.1 RBAC authentication \& authorization

Building on the initial user-service skeleton (basic JWT login with a single admin/user role) created by Lin Chuanbin, I extended the platform's security model into a multi-layer permission architecture:

- Designed five RBAC tables supporting three role types (admin, department admin, regular user) with fine-grained permissions pre-seeded in the database.
- Replaced the original single-role JWT filter with a custom JwtAuthenticationFilter that extracts roles and permissions from JWT claims and registers them in Spring Security, enabling declarative @PreAuthorize checks.
- Built admin and department-admin controllers covering user CRUD, role assignment, department assignment, clearance level management, and account lifecycle operations.
- Introduced multi-level department trees with automatic department-scoped data isolation in all data-access controllers.
- Implemented three-tier document clearance enforcement with department scoping, clearance level comparison, and HITL override via approved access requests.

1.2 API gateway, session management \& unified response format

I built the platform's unified API gateway layer and standardized the cross-service response format:

- Built a custom Spring Cloud Gateway global filter with route rules covering all microservices and WebSocket endpoints.
- Implemented single-session enforcement via Redis (new logins invalidate old sessions), sliding-window TTL renewal with polling endpoint exemption, and dynamic admin-configurable session timeout.
- Propagated user identity through HTTP headers so downstream services operate without holding JWT\_SECRET.
- Defined the Result\textless T\textgreater generic wrapper standard, deployed identical DTOs across all four microservices, and wrote the frontend Axios interceptor with automatic unwrapping.

1.3 Notification \& HITL approval system

I independently designed and implemented the platform's notification and approval subsystem:

- Designed the notification table with threaded messaging support (parent\_id/thread\_id) and a five-state state machine spanning unread, read, pending, approved, and rejected statuses.
- Built the notification controller with endpoints for sending, filtered listing, read marking, approve/deny with audit trail, and threaded conversation view.

1.4 Task center orchestration

I designed and implemented the task center, a centralized orchestration layer that tracks and manages all Agent workflows across the platform. When a user submits a request through any Agent, the task center persists it as a task record, queues it for asynchronous execution, and pushes real-time progress updates to the frontend.

- Built the task-service microservice with an async execution pipeline.
- Implemented routing by task type to the corresponding Agent, with progress pushed in three stages to the frontend.
- Designed fault-tolerance mechanisms including connection timeout, read timeout, up to 3 automatic retries, and graceful queue overflow rejection to prevent system overload.

1.5 Email verification system

I replaced the initial QQ SMTP implementation (originally contributed by Xia Xin) with a production-grade solution after identifying reliability and security limitations:

- Migrated from QQ SMTP to the Resend HTTP API with a redesigned HTML email template.
- Hardened security with brute-force protection against verification code guessing and a three-tier rate limiting system (per-email cooldown, per-IP hourly cap, per-email daily cap).

1.6 Copilot backend \& AI provider unification

I built the backend for the global Copilot assistant, which generalized the AI chat panel originally confined to Lin Chuanbin's Tool Agent page into a cross-page, multi-Agent intelligent dispatcher that sits above all three agents as a unified natural language entry point. The Copilot's Python Flask intent classifier implements a three-tier classification pipeline --- security blocklist, keyword greeting match, and LLM-based six-category classification (CODE, TOOL, RAG, CHAT, CLARIFY, SETTINGS) --- complemented by multi-turn clarification, cross-page result delivery via browser custom events, and WebSocket-based progress streaming for long-running agent calls. I unified all AI provider configurations across the platform into a single admin-managed configuration center with AES-256 encryption, eliminating the fragmentation that existed when each agent managed its own LLM settings independently.

1.7 Document management \& MinIO storage

I worked with Huang Siyuan to build the document storage and management backend, which serves both the document library and her RAG Agent:

- Built a MinIO storage service with auto bucket creation, admin-configurable upload size limits, and support for PDF, DOCX, PPT, and Markdown files.
- Extended the document table with file metadata and parse status columns to support RAG document parsing.
- Implemented document access request approval routing to department managers with fallback to system admin.

\textbf{2. Frontend}

2.1 Layout, navigation \& multi-tab system

I rebuilt the frontend shell from the initial scaffold created by Lin Chuanbin into a production-grade application framework:

- Rebuilt the main layout with a collapsible sidebar, multi-tab top bar, responsive mobile layout with off-screen drawer, and role-based dynamic menu visibility.
- Built a browser-style multi-tab system that preserves page state when users switch between tabs --- form inputs, filters, and scroll positions remain intact --- preventing task interruption, a common issue in enterprise workflows.
- Implemented a five-layer global navigation guard covering silent refresh recovery, authentication redirect, bounce-back, single-role check, and multi-role check.

2.2 UI design system

I established a minimalist visual design language across the entire platform, replacing the default Element Plus styling inherited from the initial scaffold:

- Globally replaced the default blue theme with a grayscale color system, removing accent strips and unifying border-radius, sizing, and spacing.
- Redesigned notifications with macOS-style frosted glass and custom cubic-bezier slide-in animation.

2.3 Tool Agent UI rebuild

I redesigned all three Tool Agent interfaces originally built by Lin Chuanbin, splitting the monolithic ~600-line file into independent sub-components and removing the fixed AI chat panel in favor of an inline interaction model:

- Meeting Agent: rebuilt search as a horizontal inline bar with a responsive card grid, SVG icons, equipment tags, and a booking dialog.
- Schedule Agent: unified two separate forms into a single filter panel and replaced manual user ID input with a filterable multi-select dropdown.
- Route Agent: redesigned the map from conditional rendering (hidden when no results) to a persistent full-width display. Added a floating control panel and clickable route segments that show detailed navigation steps.

\textbf{3. DevOps}

I set up the full deployment pipeline for the platform, bringing all three Agent services --- Zhang Zhaoming's Code Agent, Huang Siyuan's RAG Agent, and Lin Chuanbin's initial Tool Agent --- onto a single DigitalOcean Droplet:

- Wrote the docker-compose.yml that starts all 9 services with proper startup order, health checks, and data persistence.
- Created Dockerfiles for each Spring Boot service.
- Deployed a DigitalOcean Droplet. To keep the Maven build from running out of memory on a small instance, I compiled the JARs locally and copied them over.
- Set up Nginx as the only public-facing entry point, handling SPA routing, API/WebSocket proxying, and static file caching.
- Registered the domain bankagent.online through Namecheap and configured Let's Encrypt for HTTPS access.

\vspace{4mm}

\noindent\textbf{Huang Siyuan:}

\vspace{2mm}

I was responsible for building the RAG Agent into a complete enterprise document Q\&A subsystem. I created a standalone rag-agent microservice on port 8085, integrated it with the Spring Cloud Gateway via the new /api/rag/** route, designed three new database tables (rag\_document\_chunk, rag\_milvus\_vector, rag\_query\_log) and integrated Milvus as the vector database with HNSW/IVF\_FLAT index for approximate nearest neighbor retrieval. I implemented document chunking with structure-aware splitting and configurable overlap, embedding via external API and local ONNX models, and permission-enforced knowledge retrieval through pre-retrieval scope computation and post-retrieval re-verification. I defined five REST API endpoints (POST /rag/query, POST /rag/index/rebuild, POST /rag/index/document/\{id\}, DELETE /rag/index/document/\{id\}, GET /rag/health) protected by Gateway-level JWT validation. I also built the frontend RAG workbench page with a split-pane layout: left-side question input with example prompts, right-side answer panel with cited source documents, and a bottom toolbar with permission filter indicators and index management controls.

\vspace{4mm}

\noindent\textbf{Xia Xin:}

\vspace{2mm}

I contributed across backend security, frontend internationalization, documentation standardization, authentication, multilingual support, and visual communication of system architecture.

\textbf{1. Backend}

1.1 Authentication

I built the platform's initial Spring Security layer with stateless JWT tokens (HMAC-SHA256) and BCrypt password encoding to separate admin and regular user access. I configured the security filter chain to expose only the login and registration endpoints publicly while requiring authentication for all other routes, and defined the foundational Result<T> response wrapper and DTOs to ensure consistent error handling and data exchange across the codebase. This part was upgraded by Yao Jiacheng.

1.2 Email Verification Integration
Building on the existing team codebase, I integrated the QQ Mail SMTP service to deliver verification codes during registration. I stored codes in Redis with a 5-minute TTL, enforced a 60-second per-email send rate limit to prevent abuse, and wired the code verification step directly into the registration flow, completing the end-to-end sign-up process. This part has been upgraded into Resend HTTP API by Yao Jiacheng.

\textbf{2. Frontend & Internationalization}

2.1 Multi-Language i18n Support
I adapted all business and UI components to support three languages (Simplified Chinese, Traditional Chinese, and English) with dynamic switching through Vue3 and Element Plus i18n. This required systematically reorganizing component-level text, validation messages, and document metadata into structured locale files, ensuring the entire frontend remained maintainable and language-agnostic. Also, I designed a initial version of Code Agent's frontend, which has been upgraded into a better version by Yao Jiacheng.

\textbf{3. Documentation & Visualization Standardization}

3.1 System Diagrams Redesign
I designed the first version of all figures, and after team discussion, I redrew the original dense, small-font diagrams into clean, wide-format vector graphics with a unified visual style. Every diagram maintained readability and visual hierarchy, ensuring clear communication of complex architectural concepts. A total of 12 figures were created and refined.

3.2 Final Report & Table Standardization

I was responsible for synchronizing the final report with the latest codebase. I converted all seven data tables into the three-line table format (1.5 pt top and bottom rules, a 0.75 pt rule below the header, no vertical or interior horizontal lines) and suggested moving long pseudocode blocks to standalone tables—both to maintain formatting consistency and to communicate algorithms more clearly to readers.

\vspace{4mm}

\noindent\textbf{Zhang Zhaoming:}

\vspace{2mm}

I designed and implemented the Code Agent module — the natural-language-to-SQL engine that translates user questions into MySQL queries, validates them through a five-layer security whitelist, and executes them safely against the database. The module comprises a dedicated Spring Boot microservice and a Python LLM inference bridge. I also conducted model-level experimentation, fine-tuning two T5 variants before converging on the final LLM API-based approach.

\textbf{1. Code Agent Microservice Architecture}

I built the code-agent microservice (port 8084) from scratch as a standalone Spring Boot 3.2 service.

Designed the service with a clean separation of concerns: CodeAgentController (REST layer, 6 endpoints), CodeGenerationService (interface with pluggable implementations), SqlValidationService (five-layer whitelist engine), SqlExecutionService (JdbcTemplate-based safe execution), MetadataCacheService (information_schema → Redis), and MetadataCacheManager (lifecycle management).

Implemented the OnnxCodeGenerationService as an HTTP proxy to the Python LLM inference server (port 8090), using java.net.http.HttpClient with 10-second connect timeout and 60-second response timeout. The contract is POST {question} → returns {sql, method} JSON.

Built the CodeAgentProperties configuration class with @ConfigurationProperties(prefix = "code-agent"), supporting three config groups (ONNX inference, SQL whitelist, metadata cache) all overridable via environment variables for seamless Dev/Prod switching.

\textbf{2. Five-Layer SQL Whitelist Validation Engine}

I designed and implemented a defense-in-depth SQL validation pipeline (SqlValidationService) that blocks injection and unauthorized queries at five sequential layers, each with immediate rejection on failure:

Layer 1 — Operation Type: Enforces SELECT-only. Any non-SELECT statement is rejected before further processing, preventing write operations at the earliest stage.

Layer 2 — Forbidden Keywords: Scans SQL with word-boundary regular expressions (\bDROP\b, \bDELETE\b, \bINSERT\b, \bUPDATE\b, \bALTER\b, \bTRUNCATE\b, \bCREATE\b, \bEXEC\b, \bUNION\b, \bINTO\b, \bLOAD_FILE\b, \bOUTFILE\b, \bDUMPFILE\b). The \b boundary prevents false positives — for example, customer_no is not flagged for containing the substring IN.

Layer 3 — Table Whitelist: Extracts all tables referenced in FROM and JOIN clauses via regex and validates each against the cached information_schema table list. Unknown tables are rejected with the allowed set shown in the error message.

Layer 4 — Column Whitelist: For non-SELECT * queries, extracts every column from the SELECT clause and validates each against the actual column definitions of the referenced tables — supporting JOIN queries by merging all referenced tables' column sets. Aggregate functions and expressions are intelligently skipped.

Layer 5 — Complexity Limits: Caps queries at a maximum of 3 joined tables and 10 conditions (counting AND/OR occurrences + 1 base). This prevents resource exhaustion from overly complex LLM-generated queries.

All layer parameters are configurable through code-agent.whitelist.* in application.yml.

\textbf{3. Metadata Caching & Execution Infrastructure}

I built the metadata caching layer and safe SQL execution service:

Implemented MetadataCacheService that queries MySQL information_schema (using JDBC directly, not MyBatis) to extract table names, column names, data types, and comments — caching them in Redis as a JSON Hash map keyed by table name with a 1-hour TTL. When Redis is unavailable, the service falls back to direct MySQL queries, ensuring zero-downtime operation.

Built MetadataCacheManager with an @EventListener(ApplicationReadyEvent.class) for startup cache warm-up and a @Scheduled(fixedRate = 30 * 60 * 1000) for periodic refresh every 30 minutes.

Implemented SqlExecutionService with defense-in-depth: before executing any SQL, it re-validates through the full five-layer pipeline, catching edge cases where SQL might have been tampered with between generation and execution. Only then does it run via JdbcTemplate.queryForList(sql) and return results with column names and data rows.

Built the generateAndExecute one-shot method that chains generation → validation → execution into a single FullResult response.

\textbf{4. REST API Endpoints}

I designed and implemented all six REST endpoints for the Code Agent, centered around the Human-in-the-Loop (HITL) pattern that separates generation from execution:

POST /code/query is the one-shot endpoint: it accepts a natural language question, calls the LLM inference service to generate SQL, runs the five-layer validation, executes the query, and returns the full result including column headers, data rows, row count, and elapsed time.

POST /code/generate handles SQL generation only — it takes a question and returns the LLM-generated SQL with inference metadata, without executing anything. This enables the frontend to display the generated SQL for human review before it touches the database.

POST /code/execute accepts a reviewed (and potentially human-edited) SQL string, re-validates it through the full whitelist pipeline, and executes it — completing the supervised HITL workflow.

GET /code/metadata returns the currently cached table metadata from Redis or MySQL fallback.

POST /code/metadata/refresh forces an immediate refresh of the metadata cache from information_schema.

GET /code/health returns the service status and active inference method for operational monitoring.

\textbf{5. Model-Level Experimentation: Fine-Tuned T5 vs. LLM API}

Before settling on the final architecture, I conducted practical experiments comparing local fine-tuned models against cloud LLM APIs for the text-to-SQL task:

Fine-tuned two T5 variants on our banking-domain dataset, deploying them via ONNX Runtime for local inference. The goal was to evaluate whether a small, domain-specific model could match general-purpose LLMs on our constrained schema.

Through systematic comparison, I found that the fine-tuned T5 models consistently underperformed the LLM API approach — even with careful training, the smaller models struggled with schema grounding, complex JOIN logic, and edge cases that a general-purpose LLM handled naturally when given well-structured metadata in the prompt.

This led to the key architectural decision: rather than investing further in local model optimization, I adopted the LLM API + rich metadata prompt strategy, where the cached information_schema context (table names, column definitions, data types, example values) is injected into every generation request. This approach proved both more accurate and more maintainable, as schema changes require no model retraining — only a metadata cache refresh.

\textbf{6. Test Dataset & Banking Domain}

I authored the banking_data.sql test dataset with three relational tables — bank_customer (5 rows, risk levels LOW/MEDIUM/HIGH), bank_account (7 rows, account types SAVINGS/CHECKING/FIXED), and bank_transaction (10 rows, transaction types DEPOSIT/WITHDRAW/TRANSFER) — covering diverse NL2SQL scenarios from simple lookups to multi-table JOIN aggregations, and providing a realistic testbed for the generation, validation, and execution pipeline.

\end{document}
